\section{Discussion}
\label{sec:discussion}

\begin{table}[t]
\caption{Where the rule stops and what is not measured. Each row is a condition under which the
rule is expected to weaken or to be the wrong tool; \emph{measured} gives the endpoint we have
for it in this paper, and a dash means we did not run that configuration. We list the
unmeasured rows on purpose: they bound what the evidence supports more than the measured ones
do.}
\label{tab:scope}
\centering\small
\setlength{\tabcolsep}{3pt}
\resizebox{\columnwidth}{!}{%
\begin{tabular}{lll}
\toprule
Condition & Current behaviour & Measured \\
\midrule
Question names no entity (topical)   & rule not needed            & $+0.0013$ (n.s.) \\
Anchor is not itself protected       & operation barely matters  & $0.6416$ vs $0.6473$ \\
Answer outside the fallback window   & refill order decides       & $0.6733$ two-hop \\
Anchor block larger than 20 lines    & block spends the budget    & $0.257$ multi \\
Anchor and answer in different blocks& routing keeps wrong block  & --- \\
Two entities share one identifier    & anchor is ambiguous        & --- \\
Alias map disclosed to the reader    & full recovery by design    & --- \\
Many prompts seen by one provider    & cross-prompt linkage       & --- \\
\bottomrule
\end{tabular}}
\end{table}


\paragraph{What the study establishes.}
For excerpts of IoT-style operational records, queried by a question that names an entity
whose value must not reach the model, the record that answers the question and the record
that stores the value are the same block. Routing each unit on (sensitivity, relevance) and
aliasing the sensitive-and-relevant cell therefore preserves the task at half the tokens and
at zero measured leakage, without a reference answer and without touching the reader. The
claim is that specific: it is a statement about queries whose referent is also the protected
value, in a layout where the two coincide.

\paragraph{Where the mechanism stops.}
The three variants in Section~\ref{sec:analysis} delimit it. The rule is unnecessary when the
question is topical and names no entity, since there is then no referent to protect or to
recover. It is redundant when anchor and answer coincide in a single record, where protecting
the anchor line already suffices. It depends on the refill order when the answer lies outside
the fallback window, the one place where the ordering rather than the routing decides the
outcome. On the topical variant, whose questions name no entity, there is neither referent to
protect nor referent to recover, and the rule has nothing to act on.

\paragraph{The adversary we assume.}
The adversary is the model provider: it receives the prompt, may retain it, and may hold many
prompts from the same deployment. It does not hold the alias map, does not observe the raw log,
and does not observe the local construction, which runs inside the operator's boundary. Under
that model the quantity we bound is the recoverability of a protected value from the prompt
text, and the routing rule reduces it to zero on this benchmark.

Three properties of the construction belong in this paragraph rather than in the evidence,
because they are choices and not measurements. Aliases are assigned per prompt, not from a
global table: the value pool in Section~\ref{sec:benchmark} is small enough that a handle
recurs in roughly twenty cases, so a global pseudonym would let one provider link those cases
to each other. Per-prompt assignment removes that link at the cost of making the map
prompt-local. The assignment is keyed, which removes the rank attack that an unkeyed
assignment admits, but a keyed assignment is still an assignment: whoever holds the key
recovers every value. And nothing here prevents the reader from \emph{inferring} a protected
attribute that the context still implies, because answering such questions is the task.

\paragraph{What the privacy claim covers.}
Leakage is measured as recoverability of the protected values from the exact prompt text,
using the pattern set of the published protocol (full value plus fragments). That is the
relevant quantity for a pipeline whose adversary is the model provider and whose prompts may
be retained.

Pattern matching answers whether a value is \emph{present}; it does not answer whether a model can
\emph{find} it, so we also run the extraction attack the same adversary would run. A model is given
each arm's passage with no question and asked to enumerate any account identifiers, addresses or
telephone numbers it can recover. On the uncompressed prompt it recovers $36.4\%$ of the protected
values and returns at least one in $82.0\%$ of cases ($[80.6,83.4]$). On \textsc{DARK} it returns
none: $0.0\%$ of cases, with a $95\%$ upper bound of $0.1\%$. The zero is not an artifact of the
pattern set; a model actively searching for the values finds nothing to look at.

It is not a formal guarantee. We derive no differential-privacy
bound~\cite{zhao-dp-2024} and no $k$-anonymity
guarantee~\cite{carvalho-microdata-2023}; the alias map is assumed private, and an adversary who
obtains it recovers the values. Keying the assignment removes the rank attack an unkeyed
assignment admits---an adversary holding only the value set can otherwise invert the map.

Both zeros are conditional on knowing which values to protect, and that condition is measurable
rather than assumed. Holding the $500$ case set fixed and keeping each protected value with
probability $\rho$, a detector that misses one value in a hundred already leaves a value
recoverable in $7.8\%$ of cases; missing one in twenty reaches $36.4\%$, and missing one in ten
reaches $58.4\%$. The task metric barely moves across that range ($0.491$ to $0.510$) while leakage
grows by more than an order of magnitude, which makes detector recall the operating requirement for
a deployment rather than a secondary detail. Below $\rho=0.9$ the rule stops protecting what it
cannot see, and the F1 gain in that regime is the arithmetic of protecting less.

The measured benefit of keeping repetitive units follows from an anti-correlation between
redundancy and answer-bearing in this data. It is not derived from the definition of the
score, and a corpus whose answers are unique rather than repeated would not exhibit it. The
principled fix is to make the relevant region cover the answer instead of relying on the
ordering. In the two-hop variant the answer lies at a fixed distance of six sentences from
the anchor, so a window of that width would protect it explicitly and remove the dependence.
We have not run that configuration, and the present results should be read as depending on
the fallback width.

Table~\ref{tab:scope} collects the conditions under which the rule weakens. Four of the eight
rows are measured; we list the other four because they bound the claim more than the measured
ones do. The two that a reader should weigh first are the cross-prompt one, which our threat
model assumes away rather than defeats, and the block-size one, where the rule does not fail
gracefully: it still returns a prompt inside budget and the loss appears only in the score.

\paragraph{Cost and threats.}
Prompt construction is lexical and model-free: one pass over the units to label them, one
pass to score redundancy (quadratic in the number of units in the worst case), and one pass
to fit the budget. It adds no model call. On the longest excerpt in the extension set
($19{,}838$ tokens uncompressed) construction completes in under a second on CPU. We did not
measure end-to-end wall-clock or memory of the reader, so the efficiency we claim concerns
prompt length rather than latency.

The protected values are synthetic canaries injected into public corpora, because those
corpora cannot carry real personal data. Canary leakage is measurable exactly, but the
measured rate is a property of this construction. Both benchmark sizes are generated by our
own scripts from shared templates, so the questions are entity-resolving by construction and
the two-hop distance is fixed rather than sampled. The two readers are 7--8B parameter
instruction-tuned models from different families. A third at twice that size answers
\textsc{unknown} in place of the alias on most multi-entity cases, where every arm then ties on
identical prompts, so the paper reports the two readers on which the mechanism operates. The two
benchmark sizes reuse case identifiers for different content, so a join across them by identifier
would silently mix records; the per-set reference value in the uncompressed arm catches that.
